\chapter*{Závěr}
\addcontentsline{toc}{chapter}{Závěr}

Cílem maturitního projektu bylo navrhnout a~implementovat mobilní aplikaci
LifeMeter umožňující centralizované sledování fyzické aktivity, výživy a~spánku.
Všechny požadavky vymezené v~zadání byly splněny; výsledná aplikace v~několika
ohledech zadání překračuje. Vedle primárně požadované platformy Android~10+
je aplikace funkční rovněž na iOS díky sdílené kódové základně v~jazyce
TypeScript. Integrace se systémovými zdravotními platformami byla realizována
souběžně pro obě prostředí -- prostřednictvím rozhraní Health Connect na
Androidu a~Apple HealthKit na iOS. Serverová část, jejíž existence ze zadání
přímo nevyplývala, byla vyvinuta jako plnohodnotný systém zahrnující REST API
s~52 explicitně definovanými endpointy, veřejnou webovou vrstvu, dokumentaci
ve standardu OpenAPI a~průběžné nasazení prostřednictvím GitHub Actions.

Aplikace je provozována na produkční cloudové infrastruktuře Oracle Cloud
na adrese \texttt{lifemeter.fit} a~je dostupná ke stažení prostřednictvím
webového download centra ve formě verzovaných APK buildů. Stejná webová
vrstva současně poskytuje veřejnou prezentaci produktu a~administrátorskou
konzoli pouze pro čtení. Uživateli nabízí jednotné prostředí pro oblasti, které stávající řešení
na trhu pokrývají odděleně a~zpravidla za poplatek -- evidenci silových
tréninků s~výkonnostními grafy, sledování výživy s~databází přes milion
potravin ověřených zdrojem USDA FoodData Central, záznamy spánku s~importem
fází z~chytrých zařízení a~přehledovou domovskou obrazovku kombinující
všechny sledované metriky. Offline režim zajišťuje plnou funkčnost i~bez
aktivního síťového připojení.

Projekt přinesl praktickou zkušenost s~celým vývojovým cyklem
full-stack mobilního a~webového systému -- od návrhu datového modelu a~architektonických
rozhodnutí přes implementaci a~automatizované nasazení až po produkční provoz.
Zvolená architektura s~důsledným oddělením klientské, serverové a~datové vrstvy
vytváří základ, který umožňuje aplikaci dále rozvíjet v~směrech nastíněných
v~kapitole Diskuze.
